\chapter{Architecture}
\section{Tổng quan Kiến trúc (System Overview)}

Hệ thống KBMS (phiên bản V3) được thiết kế dựa trên mô hình kiến trúc 4 tầng (4-Tier Architecture) hiện đại [1]. Cấu trúc này giúp tách biệt rõ ràng trách nhiệm giữa việc tương tác người dùng, truyền dẫn mạng, xử lý logic và lưu trữ vật lý [3].

\subsection{Mô hình Kiến trúc 4 Tầng}

Cấu trúc phân tầng của KBMS bao gồm:

\begin{figure}[ht!]
\centering
\includegraphics[width=0.8\textwidth]{assets/kbms\_4\_tier\_architecture.png}
\caption{Kiến trúc 4 tầng chuẩn hóa của hệ thống KBMS V3.}
\label{fig:4_1}
\end{figure}

1.  \textbf{Application Layer}: Giao diện người dùng (Studio/CLI), chịu trách nhiệm thu thập yêu cầu và hiển thị kết quả.
2.  \textbf{Network Layer}: Giao thức truyền dẫn nhị phân, đảm bảo kết nối ổn định và hiệu năng cao.
3.  \textbf{Server Layer}: "Bộ não" của hệ thống, thực hiện biên dịch ngôn ngữ KBQL và suy diễn tri thức.
4.  \textbf{Storage Layer}: Quản lý dữ liệu bền vững trên đĩa cứng thông qua các thuật toán B+ Tree và WAL.

\subsection{Luồng Yêu cầu Tổng quát}

Mọi yêu cầu từ người dùng đều đi xuyên suốt qua 4 tầng này để đạt được kết quả cuối cùng:

\begin{figure}[ht!]
\centering
\includegraphics[width=0.8\textwidth]{assets/kbms\_request\_flow\_v3.png}
\caption{Luồng xử lý yêu cầu đi xuyên suốt 4 tầng kiến trúc.}
\label{fig:4_2}
\end{figure}

\subsection{Bảo mật & Chẩn đoán Hệ thống}

Bên cạnh luồng dữ liệu chính, KBMS duy trì một "mạch quản trị" song song để đảm bảo tính an toàn và minh bạch:

\begin{figure}[ht!]
\centering
\includegraphics[width=0.8\textwidth]{assets/kbms\_security\_diagnostics\_flow.png}
\caption{Sơ đồ luồng bảo mật và chẩn đoán hệ thống song song.}
\label{fig:4_3}
\end{figure}

\subsection{Công nghệ Sử dụng (Technology Stack)}

| Tầng | Thành phần chính | Công nghệ \& Thư viện |
| :--- | :--- | :--- |
| \textbf{Application} | Studio IDE | React, Electron, Monaco Editor, Tailwind CSS |
| \textbf{Network} | Binary Protocol | TCP Sockets, UTF-8 Encoding, Little-Endian |
| \textbf{Server} | Core Engine | .NET Core 8, C\#, F-Closure Algorithm |
| \textbf{Storage} | Physical Layer | B+ Tree, WAL, Slotted Page, AES-256 Encryption |

---

> [!TIP]
> \textbf{Khả năng mở rộng}: Với kiến trúc 4 tầng chuẩn hóa, KBMS dễ dàng hỗ trợ các ứng dụng bên thứ ba tích hợp thông qua Giao thức Mạng mà không cần can thiệp vào mã nguồn Server.

\section{Tầng Ứng dụng (Application Layer)}

Tầng Ứng dụng là điểm chạm đầu tiên giữa người dùng và hệ thống KBMS. Nó bao gồm hai thành phần chính phục vụ các mục đích khác nhau.

\subsection{KBMS Studio (IDE Chuyên nghiệp)}

KBMS Studio được xây dựng trên nền tảng React và Electron, đóng vai trò là một môi trường phát triển tri thức toàn diện (Knowledge IDE).

*   \textbf{Chức năng chính}:
    *   \textbf{Monaco Engine}: Soạn thảo mã nguồn KBQL với tính năng IntelliSense và báo lỗi trực tiếp.
    *   \textbf{Visual Knowledge Designer}: Cho phép thiết kế mô hình tri thức bằng đồ thị kéo thả.
    *   \textbf{Management Dashboard}: Giám sát hiệu năng Server, quản lý User và xem Audit Logs.
*   \textbf{Công nghệ}: React, TypeScript, Monaco Editor, Tailwind CSS.

\subsection{KBMS CLI (Command Line Interface)}

KBMS CLI là công cụ dòng lệnh mạnh mẽ dành cho các nhà quản trị và nhà phát triển hệ thống.

*   \textbf{Chức năng chính}:
    *   \textbf{REPL (Read-Eval-Print Loop)}: Môi trường tương tác dòng lệnh trực tiếp với Server.
    *   \textbf{Batch Processing}: Thực thi các kịch bản tri thức (.kbql) quy mô lớn thông qua lệnh `SOURCE`.
    *   \textbf{Advanced Formatting}: Hỗ trợ hiển thị kết quả dưới dạng bảng hoặc hàng dọc (`\G`).
*   \textbf{Công nghệ}: .NET Core, C\#, `LineEditor` custom library.

---

> [!TIP]
> \textbf{Lựa chọn công cụ}: Studio phù hợp cho công tác nghiên cứu và thiết kế mô hình; CLI phù hợp cho công tác bảo trì và tự động hóa hệ thống.

\section{Tầng Mạng (Network Layer)}

Tầng Mạng đóng vai trò là "mạch máu" truyền dẫn thông tin giữa các thành phần Application Layer và Server Layer.

\subsection{Giao thức Nhị phân (Custom Binary Protocol)}

Để tối ưu hóa hiệu năng và giảm thiểu băng thông, KBMS không sử dụng JSON/XML để giao tiếp mà sử dụng một giao thức nhị phân tùy chỉnh trên nền TCP.

*   \textbf{Cấu trúc Gói tin}: Gồm 1 byte `MessageType` và 4 byte `PayloadLength`, sau đó là `Payload` dữ liệu thực tế.
*   \textbf{12+ Message Types}: 
    *   `LOGIN`: Xác thực người dùng.
    *   `QUERY`: Gửi yêu cầu KBQL.
    *   `LOGS\_STREAM`: Yêu cầu streaming log thời gian thực.
    *   `MANAGEMENT\_CMD`: Lệnh quản trị hệ thống (Dành cho ROOT).
*   \textbf{Encoding}: Sử dụng chuẩn UTF-8 cho chuỗi văn bản và định dạng nhị phân Little-Endian cho các số nguyên/thực.

\subsection{Quản lý Kết nối (Connection Manager)}

*   \textbf{Asynchronous Socket I/O}: Server sử dụng mô hình lập trình bất đối xứng để phục vụ hàng trăm kết nối đồng thời với độ trễ tối thiểu.
*   \textbf{Session Isolation}: Mỗi kết nối được gán một `SessionId` duy nhất, giúp tách biệt dữ liệu tạm thời và ngữ cảnh làm việc (Context) giữa các người dùng.
*   \textbf{Heartbeat \& Timeout}: Đảm bảo giải phóng tài nguyên hệ thống (RAM/Socket) ngay khi kết nối bị ngắt đột ngột sau 30 giây.

---

> [!IMPORTANT]
> \textbf{Hiệu năng truyền tin}: So với giao thức REST/JSON truyền thống, Giao thức Nhị phân của KBMS giúp giảm tới 60\% kích thước gói tin và 40\% thời gian xử lý tuần tự hóa (Serialization).

\section{Tầng Máy chủ (Server Layer)}

Tầng Máy chủ (Server Layer) là thành phần phức tạp nhất của KBMS, chịu trách nhiệm cho toàn bộ quá trình biên dịch ngôn ngữ và suy diễn tự động.

\subsection{Knowledge Manager (Orchestrator)}

Knowledge Manager là bộ điều phối trung tâm của Server. Nó chịu trách nhiệm:
*   \textbf{KB Isolation}: Sử dụng `StoragePool` để đảm bảo mỗi người dùng có thể làm việc trên một Cơ sở tri thức độc lập mà không gây xung đột tài nguyên.
*   \textbf{Thread-safe Operations}: Quản lý các lệnh nạp/hủy KB thông qua cơ chế khóa (Mutual Exclusion) để duy trì tính toàn vẹn của danh lục hệ thống.
*   \textbf{Execution Routing}: Điều phối luồng làm việc giữa Parser (Biên dịch) và Inference Engine (Suy diễn).

\subsection{KBQL Parser & Compiler}

*   \textbf{Phân tích Cú pháp}: Sử dụng thuật toán đệ quy xuống (Recursive Descent) để biên dịch câu lệnh KBQL thành cấu trúc AST (Abstract Syntax Tree).
*   \textbf{Hỗ trợ}: Biên dịch các tập lệnh KDL (Định nghĩa), KQL (Truy vấn), KCL (Bảo mật), KML (Bảo trì) và TCL (Giao dịch).

\subsection{Reasoning Engine (Bộ máy Suy diễn)}

*   \textbf{Forward Chaining}: Thuật toán suy diễn tiến dựa trên điểm đóng \textbf{F-Closure}.
*   \textbf{Mô hình COKB}: Thực hiện suy diễn trên các đối tượng tính toán phức tạp, giải quyết các bài toán về Quan hệ (IS-A, PART-OF) và Phương trình toán học.

\subsection{Dịch vụ Bảo mật & Xác thực (Security & Auth)}

Hệ thống bảo vệ tri thức thông qua cơ chế xác thực đa tầng và phân quyền dựa trên vai trò (RBAC):
*   \textbf{Authentication}: Đối soát thông tin đăng nhập với `UserCatalog` lưu trong System KB. Mật khẩu được mã hóa an toàn.
*   \textbf{Authorization (RBAC)}: Phân chia 3 cấp độ truy cập:
    *   \textbf{ROOT}: Quyền cao nhất, quản lý toàn bộ hệ thống, user và cấu hình.
    *   \textbf{ADMIN}: Quản lý các Cơ sở tri thức (KB) và Concept.
    *   \textbf{RESEARCHER}: Quyền đọc và thực thi suy diễn (Read-only / Solve).

\subsection{Hệ thống Nhật ký & Chẩn đoán (Native Logging)}

KBMS V3 sở hữu cơ chế nhật ký độc đáo, biến log thành một dạng tri thức có thể truy vấn:
*   \textbf{System Logs}: Lưu trữ các sự kiện vận hành, lỗi hệ thống và hiệu năng RAM/CPU.
*   \textbf{Audit Logs}: Ghi lại mọi câu lệnh KBQL của người dùng, thời gian thực hiện và trạng thái thành công/thất bại.
*   \textbf{Persistence}: Mọi log đều được `SystemLogger` chèn trực tiếp vào các concept `system\_logs` và `audit\_logs` bên trong \textbf{System KB}.
*   \textbf{Real-time Streaming}: Hỗ trợ đẩy log trực tiếp tới Studio thông qua WebSockets/TCP phục vụ giám sát thời gian thực.

---

> [!NOTE]
> Server Layer được thiết kế theo hướng \textbf{Thread-safe}, cho phép nhiều người dùng cùng truy cập và thực thi suy diễn đồng thời trên các vùng tri thức khác nhau.

\section{Tầng Lưu trữ (Storage Layer)}

Tầng Lưu trữ chịu trách nhiệm đảm bảo tính bền vững (Durability), an toàn (Security) và điều phối truy cập đồng thời cho tri thức.

\subsection{Ánh xạ Thành phần Hệ thống}

Hệ thống quản lý việc lưu trữ qua một phân cấp các Manager chuyên biệt:

\textbf{[Missing Image: storage_layer_map.png]}\\ \textit{Hình 4.4: Bản đồ giải phẫu thành phần tầng lưu trữ vật lý.}

*   \textbf{StoragePool (Singleton)}: Quản lý danh sách các KBs đang nạp và đảm bảo Thread-safety khi truy cập đa luồng.
*   \textbf{BufferPoolManager (LRU Cache)}: Điều phối bộ nhớ đệm RAM theo từng trang (Page).
*   \textbf{DiskManager (Physical Handler)}: Giao tiếp trực tiếp với hệ điều hành để đọc/ghi tệp nạp theo khối (16KB).
*   \textbf{WalManager (Transaction Logger)}: Quản lý tệp nhật ký `.log` để đảm bảo ACID.

\subsection{Quản lý Trang (Physical Paging)}

*   \textbf{Slotted Page Layout}: Sử dụng cấu trúc trang 16KB với \textbf{24 Bytes Header} (PageId, LSN, Prev, Next, FreePtr, Count).
*   \textbf{Hệ thống Slot}: Cho phép quản lý các bản ghi có độ dài thay đổi cực kỳ hiệu quả mà không làm phân mảnh đĩa.

\subsection{Cấu trúc Chỉ mục (B+ Tree Indexing)}

*   \textbf{Hiệu năng}: Tìm kiếm và truy xuất bản ghi với độ phức tạp $O(\log n)$.
*   \textbf{Clustered Index}: Tích hợp dữ liệu trực tiếp vào các nút lá (Leaf Nodes) của cây B+ Tree để tối thiểu hóa số lần đọc đĩa.

\subsection{Điều phối Truy cập Đồng thời (Latching & Pinning)}

Hệ thống hỗ trợ đa người dùng thao tác cùng lúc thông qua cơ chế Page Latching:
*   \textbf{Page Pinning}: Khi một Query cần dữ liệu, Bpm sẽ cung cấp một Frame trong RAM và tăng `PinCount`. Trang có `PinCount > 0` sẽ bị khóa vĩnh viễn trong RAM (Pinned), thuật toán LRU không thể loại bỏ nó cho tới khi được `Unpin`.
*   \textbf{LRU Eviction}: Khi Cache đầy, các trang có `PinCount = 0` và ít được sử dụng nhất sẽ bị trục xuất (Evicted) sau khi đã được đẩy (Flush) xuống đĩa nếu là trang bẩn (Dirty).
*   \textbf{Write-Ahead Logging (WAL)}: Mọi thay đổi dữ liệu đều được ghi nhật ký trước khi thực hiện trên RAM, đảm bảo khả năng phục hồi (Recovery) bất chấp mọi sự cố nguồn điện.

\subsection{Bảo mật Dữ liệu (Encryption)}

*   \textbf{AES-256}: Mã hóa toàn diện dữ liệu tĩnh (At-rest). Tệp dữ liệu thô trên đĩa hoàn toàn không thể đọc được nếu không có Master Key của hệ thống.

---

> [!IMPORTANT]
> \textbf{Dữ liệu \& Tri thức}: Tầng lưu trữ tách biệt tệp dữ liệu thực thể (.dat) và tệp định nghĩa tri thức (.kbf) để tối ưu hóa quá trình nạp và truy vấn.

\section{Đối soát Yêu cầu Kỹ thuật (Requirements Spec)}

Tài liệu này tổng hợp toàn bộ các yêu cầu Chức năng (Functional) và Phi chức năng (Non-functional) dựa trên kiến trúc 4 tầng của KBMS.

\subsection{Yêu cầu Chức năng (Functional Requirements)}

Hệ thống cung cấp đầy đủ các nhóm chức năng từ mức ứng dụng đến mức lưu trữ:

| Tầng Kiến trúc | Nhóm Chức năng | Hành động cụ thể |
| :--- | :--- | :--- |
| \textbf{Application} | Thiết kế \& Giám sát | IDE Monaco, Kéo thả đồ thị, Dashboad thông số, REPL nâng cao. |
| \textbf{Network} | Kết nối \& Truyền dẫn | Handshake nhị phân, Heartbeat, Giải mã Streaming Result Set. |
| \textbf{Server} | Xử lý \& Suy diễn | Biên dịch KBQL, Suy diễn F-Closure, Giải phương trình, RBAC. |
| \textbf{Storage} | Lưu trữ \& Phục hồi | B+ Tree Search, WAL Recovery, Checkpoint, AES Encryption. |

\subsection{Yêu cầu Phi chức năng (Non-functional Requirements)}

Đảm bảo hệ thống đạt được các tiêu chuẩn kỹ thuật chuyên nghiệp:

\subsubsection{Độ bền vững & Nhất quán (ACID Compliance)}
*   \textbf{Atomicity}: Mọi thay đổi dữ liệu là nguyên tố thông qua nhật ký WAL.
*   \textbf{Consistency}: Sau mọi lệnh KDL/KQL, tri thức phải ở trạng thái nhất quán về mặt logic.
*   \textbf{Durability}: Dữ liệu đã cam kết (Commit) không bị mất mát sau sự cố.

\subsubsection{Hiệu năng & Khả năng đáp ứng (Performance)}
*   \textbf{Độ trễ}: Thời gian phản hồi trung bình < 10ms trên mạng LAN.
*   \textbf{Tải}: Hỗ trợ đồng thời 256+ kết nối mà không làm giảm hiệu suất suy diễn.
*   \textbf{Truy xuất}: Tìm kiếm bản ghi $O(\log n)$ trên đĩa cứng quy mô TB.

\subsubsection{Bảo mật (Security)}
*   \textbf{Xác thực}: Cơ chế LOGIN nhị phân với mật khẩu băm (Hashed).
*   \textbf{Phân quyền}: RBAC đa tầng (Role-Based Access Control) đến từng Cơ sở tri thức.
*   \textbf{Mã hóa}: AES-256 cho dữ liệu tĩnh trên đĩa cứng.

\subsubsection{Khả năng Giám sát (Observability)}
*   \textbf{Logging}: Tự động ghi nhật ký mọi hành vi (Audit) và lỗi hệ thống (System Log).
*   \textbf{Real-time}: Stream log từ Server về Studio với độ trễ < 1s.
*   \textbf{Queryable Logs}: Cho phép dùng KBQL để truy vấn ngược lại lịch sử vận hành.

---

> [!IMPORTANT]
> Toàn bộ các yêu cầu trên được thiết kế để KBMS có thể vận hành như một hệ quản trị tri thức cấp doanh nghiệp ổn định và tin cậy.

